\documentclass{article}

\usepackage{amsmath,amssymb}
\usepackage{xurl}
\usepackage[hidelinks]{hyperref}

% Shared commands for notes in docs/paper-gaps/.

\DeclareUnicodeCharacter{2080}{\ifmmode{}_{0}\else\textsubscript{0}\fi}
\DeclareUnicodeCharacter{2081}{\ifmmode{}_{1}\else\textsubscript{1}\fi}
\DeclareUnicodeCharacter{2082}{\ifmmode{}_{2}\else\textsubscript{2}\fi}
\DeclareUnicodeCharacter{2099}{\ifmmode{}_{n}\else\textsubscript{n}\fi}

\newcommand{\Lean}{\textsc{Lean}}
\ExplSyntaxOn
\NewDocumentCommand{\leanid}{m}{
  \ifmmode
    \textsf{\detokenize{#1}}
  \else
    \group_begin:
    \urlstyle{sf}
    \tl_set:Nn \l_tmpa_tl {#1}
    \tl_replace_all:Nnn \l_tmpa_tl {\_} {_}
    \exp_args:NV \path \l_tmpa_tl
    \group_end:
  \fi
}
\ExplSyntaxOff
\providecommand{\tr}{\operatorname{tr}}
\newcommand{\tnleanIssueBase}{https://github.com/LionSR/TNLean/issues/}
\newcommand{\tnleanPrBase}{https://github.com/LionSR/TNLean/pull/}
\newcommand{\tnleanDocBase}{https://sirui-lu.com/TNLean/docs/}
\newcommand{\ghissue}[1]{\href{\tnleanIssueBase#1}{GitHub issue \##1}}
\newcommand{\ghpr}[1]{\href{\tnleanPrBase#1}{PR \##1}}
\newcommand{\doclink}[1]{\href{\tnleanDocBase#1.html}{\path{#1}}}

% Dirac notation and the identity operator, as in blueprint/src/macros/common.tex.
% \providecommand keeps note-local definitions authoritative.
\usepackage{dsfont}
\providecommand{\ket}[1]{|#1\rangle}
\providecommand{\bra}[1]{\langle #1|}
\providecommand{\braket}[2]{\langle #1 | #2 \rangle}
\providecommand{\ketbra}[2]{|#1\rangle\!\langle #2|}
\providecommand{\Id}{\mathds{1}}
\providecommand{\one}{\mathds{1}}
\providecommand{\C}{\mathbb{C}}
\providecommand{\MN}[1]{M_{#1}(\C)}
\providecommand{\MD}{M_D(\C)}

% External referencing for paper sources.  The build helper writes sanitized
% auxiliary files under build/paper-aux/ and excludes duplicated source labels.
\usepackage{xr}

% Import a generated paper auxiliary file with a caller-supplied label prefix.
% The candidate paths support builds from docs/paper-gaps/, the repository root,
% and build/paper-aux/ itself.
\newcommand{\importpaperaux}[2]{%
  \IfFileExists{../../build/paper-aux/#2.aux}{%
    \externaldocument[#1]{../../build/paper-aux/#2}%
  }{%
    \IfFileExists{build/paper-aux/#2.aux}{%
      \externaldocument[#1]{build/paper-aux/#2}%
    }{%
      \IfFileExists{#2.aux}{%
        \externaldocument[#1]{#2}%
      }{}%
    }%
  }%
}

% General external reference macros
\newcommand{\paperref}[2]{\ref{#1-#2}}
\newcommand{\papereqref}[2]{(\ref{#1-#2})}

% CPSV16 source (1606.00608 / MPDO-22-12-17-2.tex) external document and shortcuts
\importpaperaux{cpsv16-}{1606.00608/MPDO-22-12-17-2-tnlean}
\newcommand{\cpsvref}[1]{\paperref{cpsv16}{#1}}
\newcommand{\cpsveqref}[1]{\papereqref{cpsv16}{#1}}

% Verdict marker: \gapnote{<kind>}{<status>}. Kind is the policy
% classification (clarification, local-correction, scope-restriction,
% unfaithful, false-source, open-gap); status is open, wip (elimination
% underway), resolved, or historical. Machine-read by the site index;
% prints a verdict line.
\newcommand{\gapnote}[2]{\par\noindent\textbf{Verdict:} #1 (#2).\par}


\providecommand{\leanid}[1]{\textsf{\detokenize{#1}}}
\providecommand{\doclink}[1]{%
    \href{https://sirui-lu.com/TNLean/docs/#1.html}{\path{#1}}}
\providecommand{\gapnote}[2]{%
    \par\noindent\textbf{Verdict:} #1 (#2).\par}

\title{Blocking and the Parent-Hamiltonian Range in CPSV16}
\author{}
\date{2026-08-31}

\begin{document}
\maketitle
\gapnote{original-lattice parent-Hamiltonian construction}{resolved for $1<d$}

\section{The two lattice statements}

Let $A$ be a matrix product state tensor of bond dimension $D$ in the
canonical form of Cirac, P\'erez-Garc\'ia, Schuch, and Verstraete. Their sharp
block-injectivity proposition supplies a basis of normal tensors
$A_1,\ldots,A_g$ and an integer $p$ such that, writing $D_j$ for the bond
dimension of $A_j$,
\begin{align}
    1\leq p\leq 3D^5,
    \qquad
    \operatorname{span}\left\{
        \bigl((A_1^{[p]})^i,\ldots,(A_g^{[p]})^i\bigr):i
    \right\}
      &=\bigoplus_{j=1}^{g}M_{D_j}(\mathbb C).
    \label{eq:cpsv16_parent_hamiltonian_range_block_span}
\end{align}
The discussion following Definition~3.9 then states both that $A$ has a
parent Hamiltonian of range at most $3D^5$ and that, after blocking, the
nearest-neighbour parent Hamiltonian has the expected ground space
\cite[lines~511--527]{Cirac2016MPDO_arXiv}.

The second statement has a direct meaning on the blocked lattice. For every
blocked-chain length $N>2$, it asserts
\begin{align}
    \ker H_2^{(N)}\!\left(A^{[p]}\right)
      &=\operatorname{span}\left\{
          V^{(N)}\!\left(A_j^{[p]}\right):1\leq j\leq g
        \right\}.
    \label{eq:cpsv16_parent_hamiltonian_range_blocked_statement}
\end{align}
This is the statement proved in the formal development.
\footnote{Formal declaration:
\leanid{MPSTensor.IsCPSVCanonicalForm.exists_bnt_blocked_hasNearestNeighborParentHamiltonian}
in \doclink{TNLean/MPS/ParentHamiltonian/CPSVBlockedNearestNeighbor}.}

The first statement concerns every sufficiently long chain on the original
lattice. It asserts that there are BNT representatives $A_1,\ldots,A_g$ and a
positive integer $L\leq3D^5$ such that
\begin{align}
    \ker H_L^{(N)}(A)
      &=\operatorname{span}\left\{V^{(N)}(A_j):1\leq j\leq g\right\}.
    \label{eq:cpsv16_parent_hamiltonian_range_original_statement}
\end{align}
For $1<d$, the formal development proves this all-chain spanning clause
simultaneously with $D^2<d^L$. Thus it also verifies that the local
orthogonal-complement interaction is nonzero.
\footnote{Formal declaration:
\leanid{MPSTensor.IsCPSVCanonicalForm.exists_bnt_hasParentHamiltonianGroundSpaceSpanning}
in \doclink{TNLean/MPS/ParentHamiltonian/CPSVOriginalRange}.}

\section{Why the blocked theorem alone is insufficient}

A two-site interaction on the blocked lattice acts on $2p$ consecutive sites
of the original lattice. Consequently,
\begin{align}
    \text{nearest-neighbour range for }A^{[p]}
      &\Longrightarrow \text{range }2p\text{ for the induced interaction on }A,
    \label{eq:cpsv16_parent_hamiltonian_range_naive_transport}
\end{align}
not range $p$. Moreover, the quantifier $N>2$ in
(\ref{eq:cpsv16_parent_hamiltonian_range_blocked_statement}) concerns original
chain lengths divisible by $p$ and greater than $2p$. Definition~3.9 instead
requires the original-lattice ground-space identity for every chain length
strictly larger than the chosen interaction range. Thus the blocked statement
alone does not establish
(\ref{eq:cpsv16_parent_hamiltonian_range_original_statement}). In
particular, transporting the two-site blocked interaction is not the proof of
the first sentence of line~527.

\section{Resolution on the original lattice}

Choose one normal representative from each BNT phase class and put each
representative in a Perron gauge such that
\begin{align}
    \sum_a A^j_a A^{j\,\dagger}_a          &=I,
    &
    \sum_a A^{j\,\dagger}_a\Lambda_j A^j_a&=\Lambda_j,
    \label{eq:cpsv16_parent_hamiltonian_range_perron_gauge}
\end{align}
where $\Lambda_j$ is positive definite. The representative dimensions obey
$\sum_jD_j\leq D$, hence $g\leq D$. Quantum Wielandt gives the common
block-injectivity length
\begin{align}
    L_0&=D^4.
    \label{eq:cpsv16_parent_hamiltonian_range_wielandt_length}
\end{align}

For $g\geq2$, the proof of
\cite[Theorem~12]{PerezGarcia2007Matrix} gives the BNT kernel equality at
range
\begin{align}
    R&=3(g-1)(L_0+1)+1
    \label{eq:cpsv16_parent_hamiltonian_range_pgvwc_range}
\end{align}
whenever $N\geq R+L_0$. Define instead
\begin{align}
    L&=3(g-1)(L_0+1)+L_0=R+L_0-1.
    \label{eq:cpsv16_parent_hamiltonian_range_absorbed}
\end{align}
Then $N>L$ implies $N\geq R+L_0$. Since $R\leq L$, antitonicity of the
cyclic local-constraint spaces and the usual frustration-free inclusion give
the sandwich
\begin{align}
    \ker H_L^{(N)}
      \subseteq \ker H_R^{(N)}
      =\operatorname{span}\{V^{(N)}(A_j):1\leq j\leq g\}
      \subseteq \ker H_L^{(N)}.
    \label{eq:cpsv16_parent_hamiltonian_range_sandwich}
\end{align}
The exact coisometric reconstruction of literal CPSV canonical form preserves
the positive-length local space, so this equality transfers from the direct
sum of representatives to $A$. Finally,
\begin{align}
    L
      &\leq3(D-1)(D^4+1)+D^4
       \leq3D^5.
    \label{eq:cpsv16_parent_hamiltonian_range_bound}
\end{align}
For one representative, one may take $L_{\mathrm{aux}}=D^4+1$. The
required periodic intersection and closure property for a normal tensor is
the result reviewed in \cite[Section~IV.C]{Cirac2021Matrix}. If there are no
representatives, one may take $L_{\mathrm{aux}}=2$; this is a formal boundary
extension of the source argument, not an additional claim extracted from
CPSV16 or PGVWC07.

To obtain one witness that also satisfies the nontriviality condition, enlarge
each auxiliary range to
\begin{align}
    L&=3D^5.
    \label{eq:cpsv16_parent_hamiltonian_range_uniform}
\end{align}
Cyclic-range antitonicity preserves the upper inclusion at the enlarged
range, and the usual frustration-free inclusion supplies the reverse one.
Therefore the spanning equality still holds for every $N>L$. If $D>0$ and
$1<d$, then
\begin{align}
    D^2<2^{2D}\leq d^{2D}\leq d^{3D^5}=d^L.
    \label{eq:cpsv16_parent_hamiltonian_range_nontriviality}
\end{align}
The same witness now satisfies $0<L$, $D^2<d^L$, and $L\leq3D^5$.

\section{The physical-dimension boundary}

The source first imposes $d^L>D^2$ and then defines the local interaction as
the orthogonal complement of the local MPS space. For $D>0$, this condition is
possible for some positive $L$ exactly when $1<d$: if $d\leq1$, then
$d^L\leq1\leq D^2$, whereas the estimate
(\ref{eq:cpsv16_parent_hamiltonian_range_nontriviality}) proves sufficiency
when $1<d$. The source-faithful formal theorem therefore states this exact
feasibility boundary explicitly and includes $D^2<d^L$ in its existential
conclusion. This is a boundary of the construction, not an extra premise used
to prove the ground-space equality. Positive ambient bond dimension is
likewise explicit through $D>0$.

\section{Conclusion}

The blocked-lattice assertion
(\ref{eq:cpsv16_parent_hamiltonian_range_blocked_statement}) is formalized
with the sharp bound $p\leq3D^5$. For $1<d$, the original-lattice assertion
(\ref{eq:cpsv16_parent_hamiltonian_range_original_statement}) is now proved
with $0<L$, $D^2<d^L$, and $L\leq3D^5$. The proof absorbs the finite-size
window of the source theorem into an auxiliary interaction range, applies the
cyclic-range sandwich, and enlarges uniformly to $L=3D^5$. This resolves the
full source construction under its necessary nontrivial physical boundary.
The blocked and original-lattice assertions remain distinct: one concerns a
nearest-neighbour interaction for the actual blocked tensor, while the other
concerns a bounded-range interaction on every sufficiently long original
periodic chain.

\bibliographystyle{alpha}
\bibliography{references}

\end{document}
